\section{Day 16: Local Properties (Nov.\ 5, 2026)}
Recall the affine communication lemma: if $P$ is a property of open affine subschemes of $X$, then
\begin{enumerate}[(i)]
    \item If $\Spec R \subset X$ satisfies $P$, then $\Spec R_f \subset X$ satisfies $P$.
    \item If $\Spec R_{f_i}$ satisfies $P$ for $(f_1, \dots, f_k) = (1)$, then $\Spec R \subset X$ satisfies $P$.
\end{enumerate}
If $X$ is covered by open affines satisfying $P$, then all open affines satisfy $P$. In this way, we can say $X$ ``is $P$''.

\begin{theorem}[Atiyah--Macdonald \S3.5] \label{thm:reduced_is_local_property}
    Being ``reduced'' is a local property, i.e., ``reduced schemes'' make sense.
\end{theorem}
\begin{proof}
    We want to show that if $R$ is reduced, then $R_f$ is reduced. If $R_{f_i}$ is reduced for $(f_1, \dots, f_k) = (1)$, then $R$ is reduced.

    First, suppose $R_f$ is \textit{not} reduced. Then there exists a nilpotent $r/f^i$, i.e., $r$ is nilpotent in $R_f$. This means $r^m = 0$ in $R_f$, so $f^N r^m = 0$ in $R$ for $m \mid N$. Thus,
    \[ \left(f^{N/m} r\right)^m = 0 \implies f^{N/m} r = 0, \]
    so $r = 0$ in $R_f$.

    Now, let $r \in R$ and suppose $r^n = 0$. Then $\restr{r^n}{D(f)} = 0$, implying $\restr{r}{D(f_i)} = 0$ for all $i$, so $r = 0$ in $R$ by the sheaf property.
\end{proof}

\begin{definition} \label{defn:integral_scheme}
    We say that $X$ is \textit{integral} if $\SO_X(U)$ is an integral domain for all open $U$.
\end{definition}
As an example, $\Spec k \times k = \Spec k \coprod \Spec k$ have functions $(1, 0)$ and $(0, 1)$ as zero divisors. Being integral is not a local property, but nonetheless, there is still a characterization for when a scheme is integral in terms of that.

\begin{theorem}[Stacks Project \S28.3] \label{thm:integral_scheme_iff_reduced_irred}
    A scheme is \textit{integral} if and only if $X$ is reduced and irreducible.
\end{theorem}
\begin{proof}
    In the forwards direction, we've already seen that $X$ is reduced, so we will now demonstrate irreducibility. Suppose $U_1$ and $U_2$ are \textit{disjoint} open subsets of $X$ (i.e., not irreducibility). The sheaf property lets us glue $1$ on $U_1$ and $0$ on $U_2$ (and vice versa) together, for which both functions multiply to $0$ in $\SO_X(U_1 \cup U_2)$.

    Conversely, suppose $a, b \in \SO_X(U)$, and let $ab = 0$. $D_U(a)$ and $D_U(b)$  have empty intersection, so without loss of generlaity, suppose $D_U(a) = \emptyset$. Consider an \textit{affine} open $\Spec R \subset U$. Then $D_{\Spec R}(a) = \emptyset$, so $a$ is nilpotent in $R$, so $a = 0$ in $R$. $U$ is covered by these, so the sheaf property yields $a = 0$, which is a contradiction.
\end{proof}

\begin{corollary} \label{cor:irreducible_scheme_has_integral_local_property}
    If $X$ is irreducible, then ``integral'' is a local property, meaning you only have to check it on an affine cover.
\end{corollary}
\begin{proof}
    In this case, $P$ is equivalent to ``reduced'', so we are done.
\end{proof}

\begin{theorem} \label{thm:noetherian_is_local_property}
    The property of being \textit{Noetherian} (i.e., every ideal is finitely genearted) is local.
\end{theorem}
We (unfortunately) call a scheme a covered by open Noetherian spectra ``locally Noetherian''. Note that Noetherian also means quasi-compact. We will do this in two steps. First, $R$ being Noetherian implies $R_f$ is Noetherian; then, $R_{f_i}$ being Noetherian for $(f_1, \dots, f_k) = (1)$ implies that $R$ is Noetherian. To do this, we will need to prove the Hilbert basis theorem.

\newpage
\begin{theorem}[Hilbert Basis Theorem] \label{thm:hilbert_basis_theorem}
    If $R$ is Noetherian, then $R[x]$ is Noetherian.
\end{theorem}
\begin{proof}
    Let us look at an ideal in $I \subset R[x]$. Choose a sequence of elements $f_1, f_2, \dots$ where $f_i$ is of minimum degree in $I \setminus (f_1, \dots, f_{i-1})$; because of this, the degrees will be weakly increasing, i.e.,
    \[ \deg f_1 \leq \deg f_2 \leq \dots. \]
    Let us look at the ideal $J \subset R$, generated by the leading coefficients $(a_i)_{i \geq 1}$ of $(f_i)_{i \geq 1}$ respectively; $R$ being Noetherian tells us that $a_{k+1} \subset (a_1, \dots, a_k)$ for some $k$, so $a_{k+1} = b_1a_1 + \dots + b_ka_k$. This means
    \[ f_{k+1} - \sum b_i x^{d_{k+1} - d_i} f_i \]
    has lower degree, which contradicts the definition of $f_{k+1}$.
\end{proof}

\begin{proof}[Proof of \ref{thm:noetherian_is_local_property}]
    Let $R_f = R[x] / (xf - 1)$. $R[x]$ is Noetherian by the HIlbert basis theorem, so we are fine.

    For the second property, let $\varphi_i \colon R \to R_{f_i}$. We want to show that, for all ideals $I \subset R$, we have
    \[ I = \bigcap \varphi_i\inv \left(\varphi_i(I) R_{f_i}\right) \]
    If we show this, then in an ascending chain of ideals,
    \[ \varphi_i\inv\left(\varphi_i(J) R_{f_i}\right) \]
    stabilizes for all $i$, so the chain stabilizes in $R$. Observe that
    \[ I \subset \bigcap \varphi_i\inv\left(\varphi_i(I) R_{f_i}\right) \]
    as $I \subset \varphi_i\inv\left(\varphi_i(I) R_{f_i}\right)$ for all $i$. We want to show the reverse inclusion: indeed, otherwise, if $r \in R$ satisfies $\varphi_i(r) \in \varphi_i(I) R_{f_i}$ for all $i$, then $r \in I$.

    We can write $r = \frac{m_i}{f_i^N}$ with $m_i \in I$. This identity holds in $R_{f_i}$. The proof of the sheaf property for rings shows that, by multiplying numerators and denominators by a big power of $f_i$, we get fractions that cancel for cross multiplying reasons. Then, use the trick that
    \[ \frac{a_1}{b_1} = \dots = \frac{a_k}{b_k} \implies \frac{a_1c_1 + \dots + a_k c_k}{b_1c_1 + \dots + b_kc_k} \]
    has denominator $1$.
\end{proof}

% \begin{definition} \label{defn:morphism_of_locally_finite_type}
%     A morphism $f \colon X \to Y$ is said to be of locally finite type if there exists a covering of $Y$ by open affine subschemes 
% \end{definition}